% =====================================================================
%  capitulos/06_anexos.tex  --  Anexos (apéndices A, B, C, ...)
%  Va después de \appendix en main.tex, por eso los \chapter se numeran A, B...
% =====================================================================

\chapter{Manual de instalación}
\label{anexo:instalacion}

EduFEM se distribuye por dos vías: un ejecutable autocontenido para uso inmediato y la instalación desde el código fuente para su inspección o extensión. El código está disponible en el repositorio público del proyecto\footnote{\url{https://github.com/Sucullani/SoftwareED}}.

\section{Ejecutable autocontenido (Windows)}

Es la vía recomendada para el alumno. Se descarga el ejecutable \texttt{.exe} publicado en el repositorio y se ejecuta directamente: no requiere instalar Python ni dependencias, ya que el empaquetado con \texttt{PyInstaller} incluye el intérprete y todas las bibliotecas. Requisito: Windows 10 u 11 de 64 bits. La generación de la memoria de cálculo en PDF necesita, además, una distribución \LaTeX{} instalada (MiKTeX o \TeX{} Live); sin ella la aplicación funciona igual y las fórmulas se rinden mediante un mecanismo alternativo.

\section{Instalación desde el código fuente}

Para docentes o desarrolladores que deseen inspeccionar o modificar el programa:

\begin{enumerate}
    \item Instalar \textbf{Python 3.10 o superior} (el desarrollo se realizó con la versión 3.13).
    \item Clonar el repositorio y ubicarse en su carpeta:
\begin{verbatim}
git clone https://github.com/Sucullani/SoftwareED.git
cd SoftwareED
\end{verbatim}
    \item Crear y activar un entorno virtual (recomendado):
\begin{verbatim}
python -m venv .venv
.venv\Scripts\activate
\end{verbatim}
    \item Instalar las dependencias:
\begin{verbatim}
pip install -r requirements.txt
\end{verbatim}
    \item Ejecutar la aplicación:
\begin{verbatim}
python main.py
\end{verbatim}
\end{enumerate}

\section{Dependencias}

El núcleo numérico y la interfaz se apoyan en bibliotecas de código abierto: NumPy y SciPy (cálculo), \texttt{ttkbootstrap} (interfaz), \texttt{pylatex} (memoria de cálculo), \texttt{PyMuPDF} (visualización del PDF), \texttt{Pillow} (figuras) y \texttt{ezdxf} (importación DXF), además de \texttt{matplotlib} y \texttt{sympy} (módulos educativos). Dos componentes son opcionales: una distribución \LaTeX{} (\texttt{pdflatex}) para la memoria de cálculo y el compilador JIT \texttt{numba} para acelerar el cálculo en mallas grandes; sin ninguno de los dos la aplicación funciona, con menor rendimiento o sin el PDF de alta calidad.


\chapter{Manual de uso}
\label{anexo:manual}

% PENDIENTE: flujo pre-proceso -> proceso -> post-proceso, atajos, ejemplos.
\pendiente{Guía de uso de la interfaz: construir el modelo, resolver, interpretar resultados.}


\chapter{Listados de código seleccionados}
\label{anexo:codigo}

% PENDIENTE: incluir fragmentos clave (motor MEF) con el paquete listings o minted.
\pendiente{Fragmentos representativos del motor de cálculo y los módulos educativos.}


\chapter{Resultados extendidos de verificación y validación}
\label{anexo:vyv}

% PENDIENTE: tablas y gráficos completos de MMS, Timoshenko y Cook.
\pendiente{Tablas completas de convergencia y validación que no entraron en el Capítulo 3.}
